% Ad Atticum 9.10 --- headnote
% ad-atticum-09-10, 18 March 49 BC, Formian villa

\textit{Cicero to Atticus, written from the Formian villa
on the fifteenth day before the Kalends of April 49 BC
(the manuscript dateline: \textit{Scr.\ in Formiano xv K.\
Apr.\ a.\ 705 (49)}). This is the day after the long
letter answering three of Atticus's together, and the day
on which --- as it turned out --- Pompey was sailing from
Brundisium for Epirus, ending the Italian war. Cicero
does not yet know it. The letter is a wakeful, semi-
private meditation, written ``with no theme proposed,'' as
if to converse with the absent friend; it is among the
most personal in the whole correspondence.}

\textit{Sections 1--3 are the meditation proper. Cicero
diagnoses himself as having been \textit{amens} from the
start in not following Pompey to the end like ``one common
soldier in the ranks'' (\textit{tamquam unus manipularis}).
He saw Pompey on 17 January, full of dread, and never
liked him after; the famous love-analogy compares
Pompey's flight to a beloved grown filthy, insipid,
indecorous, who estranges affection. ``Now the love comes
back up''; ``like that bird, I gaze out upon the sea and
long to fly away.'' Yet his hesitation was not rashness
but reasoned: he ran through the historical exempla ---
Tarquin, Coriolanus, Themistocles, Hippias, Sulla, Marius,
Cinna --- and could not see himself leading Getae,
Armenians, Colchians against Rome, or starving Italy. The
sun has fallen out of the world (the image is Atticus's
own, from some earlier letter); as long as Pompey was in
Italy, hope was alive. ``These things, these things, have
failed me.''}

\textit{Sections 4--10 are an extraordinary technical
move: Cicero unrolls the sealed volume of Atticus's
letters and walks through them in dated order, quoting
each in turn, to demonstrate to Atticus --- and to
himself --- that every step he has taken was taken on
Atticus's authority. From 10 K Febr to 7 Id. Mart, eight
letters are paraded one by one: Atticus had said Pompey
would be acting \textit{alogist\=os} if he left Italy;
that Cicero must then return to the City and not call
his stay a foreign tour; that the war was being made
\textit{aspondos}, without truce; that no patriot and
statesman (\textit{philopatris}, \textit{politikos}) should
flee with him; that flight is shameful (Pompey, Cicero
remarks, has been ``yearning-to-be-a-Sulla, yearning-to-
proscribe'' for two years past); that if Lepidus and
Volcacius remained in Italy, Cicero too should remain ---
and let himself be beaten in the contest with Pompey
rather than reign with Caesar in the coming filth. If
Lepidus and Volcacius left, Atticus had thrown up his
hands (\textit{apor\=o}); whatever happened, whatever
Cicero did, must be accepted (\textit{sterkteon}). Section
10 closes on Peducaeus's parallel approval, on Cicero's
need not for self-justification (\textit{adversus me nihil
opus est}) but for outside witnesses, and on the wakeful
consolation of having reread everything Atticus wrote
him.}
